\chapter{Glossary of Terms}
\label{app:glossary}

\begin{description}
  \item[Abelian group] A group where the operation is commutative: $a \cdot b = b \cdot a$.
  \item[Algebraic cycle] A formal linear combination of subvarieties of an algebraic variety.
  \item[Algebraic variety] The set of solutions to a system of polynomial equations.
  \item[Analytic continuation] Extending the domain of a holomorphic function beyond its original region of definition.
  \item[Boundary] The set of points of a manifold that are not interior points (the ``edge'').
  \item[Closed manifold] A compact manifold without boundary.
  \item[Cohomology] A sequence of abelian groups or vector spaces that encodes topological information about a space.
  \item[Compact] A topological space where every open cover has a finite subcover. In $\mathbb{R}^n$, equivalent to being closed and bounded.
  \item[Complex projective space $\mathbb{P}^n(\mathbb{C})$] The space of lines through the origin in $\mathbb{C}^{n+1}$, a compact complex manifold.
  \item[Conjecture] A mathematical statement that is believed to be true but has not been proven.
  \item[Elliptic curve] A smooth, projective algebraic curve of genus 1 with a distinguished rational point.
  \item[Fields Medal] The most prestigious award in mathematics, awarded every four years to mathematicians under 40.
  \item[Fundamental group] The group of homotopy classes of loops in a space, measuring its one-dimensional holes.
  \item[Gauge theory] A field theory where the Lagrangian is invariant under local transformations (gauge transformations).
  \item[Hodge decomposition] The decomposition of the cohomology of a K\"{a}hler manifold into $(p, q)$-components.
  \item[Holomorphic] Complex-differentiable; a function with a complex derivative.
  \item[Homeomorphism] A continuous bijection with continuous inverse; an equivalence of topological spaces.
  \item[Homotopy] A continuous deformation between two maps. Two loops are homotopic if one can be continuously deformed into the other.
  \item[$L$-function] A function (often a Dirichlet series) encoding arithmetic information about an algebraic object.
  \item[Manifold] A space that locally resembles Euclidean space.
  \item[Modular form] A holomorphic function on the upper half-plane with certain transformation properties under the modular group.
  \item[NP-complete] A problem in NP to which every other NP problem can be reduced in polynomial time.
  \item[Polynomial time] Running time bounded by $O(n^k)$ for some constant $k$, where $n$ is the input size.
  \item[Ricci flow] A geometric evolution equation: $\partial_t g_{ij} = -2 R_{ij}$, where $R_{ij}$ is the Ricci curvature.
  \item[Simply connected] A space with trivial fundamental group; every loop can be contracted to a point.
  \item[Singularity] A point where a function, geometric structure, or solution to an equation becomes undefined or degenerate.
  \item[Smooth function] A function with derivatives of all orders ($C^\infty$).
  \item[Topology] The study of properties of spaces preserved under continuous deformations (stretching, bending, but not tearing).
  \item[Torsion group] The subgroup of elements of finite order in an abelian group.
  \item[Turing machine] A mathematical model of computation consisting of an infinite tape, a read/write head, and a finite set of states.
  \item[Zeta function] A function that encodes arithmetic information, often expressed as a Dirichlet series. The Riemann zeta function is the prototypical example.
\end{description}
